% Corte mecánico íntegro: parte_iv.tex:220--336.
% Residencia material del ensamblaje sucesor de 102 capítulos.
% Residencia de realización: capítulo 58.
\chapter{Électron central et lecteur hélicoïdal de la matière}

\section{Unicité du centre matériel sous inversion cellulaire}

L'atlas matériel a pour domaine \(\mathcal C_{81}=\{1,\ldots,9\}^2\) et pour involution
\[
 \rho_c(r,c)=(10-r,10-c).
\]
Son unique point fixe est
\[
 (r,c)=(5,5).
\]
Les \(80\) autres cellules forment \(40\) paires. L'électron central est l'unique multisection ponctuelle ; les autres espèces internes sont des fibres, orbites ou multisections. La théorie spirale ne remet pas en question cette conclusion.

\section{Mode tritique commun}

Dans la carte canonique APP \(1,\ldots,9\), le vecteur des valeurs du sélecteur tritique de phase opératoire est
\[
 m_{\ph}=(1,-1,0)^3.
\]
Pour le tableau conjoint tridimensionnel,
\[
 C^{(3)}m_{\ph}=(C^{(3)})^{\!*}m_{\ph}
 =-\frac12m_{\ph}.
\]
Par conséquent
\[
 s=\frac12,\qquad
 1-s^2=\frac34=\frac{90}{120},
\]
\begin{equation}
 \|[P_{\Sigma,3},P_{\Pi,3}]\|
 =\frac{\sqrt3}{4}.
 \label{espiral:eq:modo-electron-app}
\end{equation}
Le mode, le défaut orienté \(90\to120\) et le préfacteur électronique sont des publications coordonnées d'une même chaîne APP--TRIT--TPK, dont la composition effective fournit leur entrelacement interne.

\section{Ledger material}

Chaque cellule de l'atlas porte la signature
\[
 \mathcal R(r,c)=
 (n_A,s_C,k_{\Delta_4},\nu_{120},\nu_{270},\tau_{12})_{r,c}.
\]
L'image contient \(56\) familles de voies ; le raffinement par famille et niveau contient \(13\) multisections. Les chiffres directeurs sont
\[
 81/56/13/324,
\]
avec le secteur nul et les tours globales. Le lecteur radial ne peut être ajouté qu'au moyen d'une signature compatible, par exemple
\[
 \Sigma_{\mathrm{esp}}(X)=
 (p_X,a_X,[\Lambda_X,C_X],
 [\kappa_{9,X}],\mu_X,\varepsilon_{{\mathrm{tw}},X}),
\]
qui soit constante là où il convient et naturelle sous raffinement.

\section{Carré de compatibilité matérielle}

Soit \(\mathcal E_{\mathrm{atlas}}\) l'espace des sections persistantes et \(\mathfrak S_{\mathrm{esp}}\) le lecteur spiral. La compatibilité matérielle s'exprime par le carré
\[
\begin{array}{ccc}
\mathcal E_{\mathrm{atlas}}
 & \xrightarrow{\ \mathfrak S_{\mathrm{esp}}\ }
 & \mathcal D_{\rad}^{\enr}\\[1.0em]
{\scriptstyle (K_D,K_\Omega)}\big\downarrow
 && \big\downarrow{\scriptstyle \ev_{\mathrm{dim}}^{\mathrm{esp}}}\\[1.0em]
\operatorname{Spec}(K_D,K_\Omega)
 & \xrightarrow{\ \mathfrak R_{\mathrm{mat}}\ }
 & \mathcal M_{\mathrm{hel}}
\end{array}
\]
L'électron \((5,5)\) fournit l'ancrage central, mais ne sélectionne pas rétrospectivement \(p,\Lambda,C\). Pour les autres espèces, on requiert la constance sur les multisections, la compatibilité avec l'inversion cellulaire, \(K_D/K_\Omega\), CKM/PMNS, la couleur et la représentation de spin.

\section{Spin et hélicité}

On n'identifie pas le spin \(1/2\) à un tour d'hélice. La représentation fermionique satisfait
\[
 U(2\pi)\psi_e=-\psi_e,
 \qquad
 U(4\pi)\psi_e=\psi_e.
\]
L'hélice nonadique enregistre la phase et la mémoire ; la représentation de spin vit dans un autre codomaine. Un opérateur d'entrelacement matériel valide doit conserver les deux lois, non remplacer l'une par l'autre.

\begin{figure}[htbp]
\centering
\begin{tikzpicture}[scale=.63]
\foreach \i in {1,...,9}{
 \foreach \j in {1,...,9}{
  \draw[HMTLine,fill=HMTLight] (\i,\j) rectangle +(1,1);
 }
}
\fill[HMTNavy] (5.15,5.15) rectangle (5.85,5.85);
\draw[HMTRed,very thick] (5.5,5.5) circle (1.05);
\draw[->,HMTRed,thick] (5.5,5.5)--(8.9,8.9);
\node[right,HMTNavy,font=\sffamily\bfseries] at (8.9,8.9) {unique point fixe \((5,5)\)};
\node[below,HMTGray,font=\small\sffamily] at (5.5,.45) {atlas material \(9\times9\)};
\end{tikzpicture}
\caption{L'électron central est un résultat de l'atlas ; la correspondance spirale est un lecteur ultérieur.}
\label{espiral:fig:atlas-electron}
\end{figure}
